\documentclass[a4paper,11pt]{article}
\usepackage[utf8]{inputenc}
\usepackage[T1]{fontenc}
\usepackage[french]{babel}
\usepackage{graphicx}
\usepackage{amsmath}
\usepackage{geometry}
\usepackage{fancyhdr}
\usepackage{hyperref}
\usepackage[usenames,dvipsnames,svgnames,table]{xcolor}
\usepackage{enumerate}
\usepackage{booktabs}
\usepackage{tcolorbox}
\tcbuselibrary{listings,skins,breakable}
\usepackage{sourcecodepro}
\usepackage{FiraSans}
\usepackage{environ}
\usepackage{tikz}
\usepackage{pgfplots}
\pgfplotsset{compat=1.18}
\usetikzlibrary{arrows.meta,positioning,calc}

\geometry{a4paper, scale=0.8}
\urlstyle{same}

%% ─── Colors ──────────────────────────────────────────────────────
\definecolor{modernblue}{RGB}{66,133,244}
\definecolor{moderngray}{RGB}{95,99,104}
\definecolor{darkgray}{RGB}{60,64,67}
\definecolor{accentorange}{RGB}{255,111,0}
\definecolor{codegreen}{RGB}{15,157,88}
\definecolor{codeblue}{RGB}{25,118,210}
\definecolor{codebg}{RGB}{250,250,250}
\definecolor{codepurple}{RGB}{156,39,176}
\definecolor{niceblue}{RGB}{0,128,172}

%% ─── Code listing: PYTHON ───────────────────────────────────────
\newtcblisting{pythoncode}[1][]{
    listing engine=listings, colback=codebg, colframe=moderngray!50,
    listing only, left=5mm, enhanced, sharp corners=all, boxrule=0.5pt, top=0pt,
    listing options={
        language=Python, basicstyle=\ttfamily\scriptsize,
        keywordstyle=\color{codeblue}\bfseries,
        commentstyle=\color{codegreen}\itshape,
        stringstyle=\color{accentorange},
        numberstyle=\tiny\color{moderngray}, numbers=left, numbersep=8pt,
        breaklines=true, tabsize=4, showstringspaces=false,
        keepspaces=true, columns=flexible,
        literate=
            {é}{{\'{e}}}1 {è}{{\`{e}}}1 {ê}{{\^{e}}}1 {ë}{{\"{e}}}1
            {û}{{\^{u}}}1 {ù}{{\`{u}}}1 {î}{{\^{i}}}1 {ï}{{\"{i}}}1
            {ô}{{\^{o}}}1 {â}{{\^{a}}}1 {à}{{\`{a}}}1 {ä}{{\"{a}}}1
            {ç}{{\c{c}}}1 {Ç}{{\c{C}}}1
            {É}{{\'{E}}}1 {È}{{\`{E}}}1 {Ê}{{\^{E}}}1,
        #1
    }, breakable
}

%% ─── Lesson box ──────────────────────────────────────────────────
\newtcolorbox{lesson}[1][]{
    enhanced, breakable,
    colback=modernblue!5, colframe=modernblue!80,
    colbacktitle=modernblue, coltitle=white,
    fonttitle=\sffamily\bfseries, fontupper=\sffamily\small,
    boxrule=0.5pt, arc=3pt, left=8pt, right=8pt, top=8pt, bottom=8pt,
    attach boxed title to top left={yshift=-2mm, xshift=5mm},
    boxed title style={sharp corners, boxrule=0.5pt},
    title={#1}, before skip=1em, after skip=2em
}

%% ─── Out-of-scope box (optional) ─────────────────────────────────
%% Use for concepts beyond the syllabus that students may encounter.
\newtcolorbox{horsprog}{
    enhanced, breakable,
    colback=codepurple!5, colframe=codepurple!60,
    colbacktitle=codepurple!80, coltitle=white,
    fonttitle=\sffamily\bfseries, fontupper=\sffamily\small,
    boxrule=0.5pt, arc=3pt, left=8pt, right=8pt, top=8pt, bottom=8pt,
    title={Hors-programme --- à lire attentivement},
    attach boxed title to top left={yshift=-2mm, xshift=5mm},
    boxed title style={sharp corners, boxrule=0.5pt},
    before skip=1em, after skip=2em
}

%% ─── Solution toggle ─────────────────────────────────────────────
\newif\ifshowsolutions
\showsolutionsfalse

\ifshowsolutions
    \newtcolorbox{solution}{colback=white, colframe=white, breakable,
        fontupper=\color{red}}
\else
    \NewEnviron{solution}{}{}
\fi

%% ─── Question command ─────────────────────────────────────────────
\newcounter{question}[section]
\newcommand{\q}[1]{%
    \vspace{0.5em}\noindent
    {\color{black}\bfseries\stepcounter{question}\Alph{section}\arabic{question}.}
    \hspace{0.5em}#1\par\vspace{0.3em}
}
\renewcommand{\thesection}{{\Alph{section}}}

%% ─── Header / footer ─────────────────────────────────────────────
\pagestyle{fancy}
\fancyhf{}
\fancyhead[L]{2025-2026}
\fancyhead[C]{\textbf{\courseName}}
\fancyhead[R]{\large GEII -- BUT2}
\fancyfoot[C]{\thepage}
\fancyfoot[L]{\includegraphics[height=3em]{../../../common/logos_couleurs/IUTA.png}}
\fancyfoot[R]{\href{mailto:prenom.nom@univ-smb.fr}{prenom.nom@univ-smb.fr}}

%% ─── Course metadata ─────────────────────────────────────────────
\newcommand{\courseName}{Outils Mathématiques et Logiciel --- OML3}
\newcommand{\chapterTitle}{TP 1 --- Série de Fourier}

\makeatletter
\renewcommand\maketitle{{\raggedright\begin{center}
    {\Huge\bfseries\@title}
\end{center}}}
\makeatother

\usepackage{amsmath,amsfonts,bm}

%% ─── Operators ───────────────────────────────────────────────────
\DeclareMathOperator*{\argmax}{arg\,max}
\DeclareMathOperator*{\argmin}{arg\,min}
\DeclareMathOperator{\MSE}{MSE}
\DeclareMathOperator{\sign}{sign}
\DeclareMathOperator{\sinc}{sinc}
\DeclareMathOperator{\tr}{tr}

%% ─── Number sets ─────────────────────────────────────────────────
\def\R{{\mathbb{R}}}
\def\C{{\mathbb{C}}}
\def\N{{\mathbb{N}}}
\def\Z{{\mathbb{Z}}}

%% ─── Infinitesimal and units ────────────────────────────────────
\def\d{{\rm d}}
\def\Hz{{\mbox{Hz}}}
\def\m{{\mbox{m}}}

%% ─── Complex / signal ───────────────────────────────────────────
\DeclareMathOperator{\complexj}{\mathsf{j}}   % imaginary unit

%% ─── Probability ─────────────────────────────────────────────────
\newcommand{\Esp}[2]{\mathbb{E}_{#1}\!\left[#2\right]}
\renewcommand{\Pr}{\mathbb{P}}

%% ─── Linear algebra ──────────────────────────────────────────────
\newcommand{\abs}[1]{\left|{#1}\right|}
\newcommand{\norm}[1]{\left\|#1\right\|}
\newcommand{\inner}[2]{\left\langle#1,\,#2\right\rangle}

%% ─── Logic ───────────────────────────────────────────────────────
\newcommand{\then}{\Rightarrow}
\newcommand{\iif}{\Leftrightarrow}

%% ─── Draft annotation ────────────────────────────────────────────
\newcommand{\red}[1]{{\color{red}#1}}

%% ─── Shared figure color (used by common/figures/*/*.tex bodies) ──
\definecolor{niceblue}{RGB}{0,128,172}

%% ─── Rappel header (styled subtitle for reference-card asides) ───
%% Requires \definecolor{niceblue}{...} to be defined by the including document.
\newcommand{\rappelhead}[1]{%
    \bigskip
    {\color{niceblue}\bfseries #1}\\[-0.4em]
    {\color{niceblue!35}\rule{\linewidth}{0.6pt}}
    \smallskip

}


\title{\chapterTitle}
\author{}
\date{}

%% ─────────────────────────────────────────────────────────────────
\begin{document}
\fontfamily{bch}\selectfont
\fontsize{12pt}{14pt}\selectfont

\maketitle
\thispagestyle{fancy}


%% ─── AVANT DE COMMENCER ──────────────────────────────────────────
\section*{Avant de commencer}

Ce TP n'est pas un TP de programmation : toutes les fonctions dont vous
aurez besoin sont déjà écrites. Votre travail consiste à exécuter les
cellules, lire les tracés produits, et répondre aux questions.

\textbf{Ce document est le sujet de référence.} Le notebook
\texttt{TP1.ipynb} est un compagnon d'exécution : il ne reprend que
l'essentiel de chaque question, repérée par son numéro (A1, A2, B1, \dots)
correspondant à celui utilisé ici. Gardez ce PDF ouvert pendant tout le TP
et répondez-y (ou sur feuille) : le notebook seul ne suffit pas.

\subsection*{Configuration de l'environnement}
\begin{enumerate}
    \item Ouvrir \texttt{TP1.ipynb} dans Jupyter.
    \item Exécuter la première cellule (import des fonctions fournies,
          dans \texttt{tp1\_helpers.py}) : rien ne devrait s'afficher, c'est
          normal.
    \item Exécuter ensuite les cellules dans l'ordre, en vous référant aux
          questions ci-dessous.
\end{enumerate}

\begin{lesson}[Un notebook, comment ça marche ?]
    Un notebook Jupyter est une suite de cellules. Une cellule de code
    s'exécute avec \texttt{Shift+Entrée} ; le résultat (tracé, valeur)
    s'affiche juste en dessous. Une cellule peut être ré-exécutée autant de
    fois que nécessaire, par exemple après avoir changé une valeur -- c'est
    ce que vous ferez pour la plupart des questions de ce TP.
\end{lesson}

\vspace{2em}\hrule\vspace{2em}
\newpage


%% ─── SECTION A ──────────────────────────────────────────────────
\section{Première manipulation de séries trigonométriques}

\subsection*{Contexte}
Avant de manipuler des séries de Fourier, on prend en main l'outil sur des
signaux déjà connus : sinus et cosinus.

\q{Afficher les fonctions sinus et cosinus entre $-10$ et $10$ (fréquence
$f=1\,\Hz$). Combien de périodes observez-vous ?}
\begin{solution}
    $T=1/f=1\,$s, donc 20 périodes sur $[-10,10]$.
\end{solution}

\q{Quel est le décalage temporel entre les deux fonctions ? Comment
reconnaître $\cos$ de $\sin$ par lecture graphique ?}
\begin{solution}
    $\cos(t)=\sin(t+\pi/2)$ : $\cos$ est en avance d'un quart de période sur
    $\sin$. Sur le graphe, $\cos$ vaut son maximum en $t=0$ tandis que
    $\sin$ y est nul (et croissant) : c'est ce qui permet de les
    distinguer visuellement.
\end{solution}

\begin{lesson}[Zoomer sans tracé interactif]
    Les tracés de ce TP ne sont pas interactifs (pas de souris pour
    zoomer). À la place, les fonctions fournies acceptent des arguments
    \texttt{xlim=(min,max)} et \texttt{ylim=(min,max)} : changez ces
    bornes et ré-exécutez la cellule pour "vous déplacer" dans le tracé.
\end{lesson}

\q{Refaire le tracé à $f=440\,\Hz$. Quelle est la période ? Zoomer sur une
période (avec \texttt{xlim}) : que devient le décalage entre $\cos$ et
$\sin$ ?}
\begin{solution}
    $T=1/440\approx 2{,}27\,$ms. Le décalage reste un quart de période
    (proportion inchangée), donc environ $0{,}57\,$ms en valeur absolue :
    seule sa durée en secondes diminue avec $f$.
\end{solution}

On fixe maintenant $T=1$ et $a_n=b_n=1$ pour tout $n$.

\begin{lesson}[Fonction \texttt{reconstruct}]
    \texttt{reconstruct(t, T, a0, an, bn)} évalue
    $a_0+\sum_n\left[a_n\cos(n\omega_0 t)+b_n\sin(n\omega_0 t)\right]$ sur
    le vecteur de temps \texttt{t}, à partir de coefficients que vous lui
    donnez (calculés ou choisis à la main).
\end{lesson}

\q{Afficher la série trigonométrique correspondante pour une valeur de $N$
de votre choix, puis recommencer avec d'autres valeurs (essayer par
exemple $N=1,2,5,20,100$, puis $N=10^5$). La période du signal change-t-elle
avec le nombre de termes ? Cette série converge-t-elle ?}
\begin{solution}
    La période ne change pas : chaque terme est de période $\le 1$, donc la
    somme reste $1$-périodique quel que soit $N$. En revanche la série ne
    converge pas : en $t=0$ chaque terme vaut $1$, donc la somme partielle
    vaut exactement $N$ -- elle diverge avec $N$ (observable numériquement :
    la valeur en $t=0$ atteint $10^5$ pour $N=10^5$). Ailleurs, la somme
    oscille sans se stabiliser. C'est attendu : rien ne garantit qu'une
    série trigonométrique quelconque converge, ce que les conditions de
    Dirichlet (vues en cours) viennent précisément garantir.
\end{solution}


%% ─── SECTION B ──────────────────────────────────────────────────
\section{Analyse et synthèse d'un signal périodique par série de Fourier}

\subsection*{Objectif}
Passer du signal carré à ses coefficients de Fourier, puis reconstruire le
signal à partir d'un nombre fini de ces coefficients, et observer ce que
cette troncature change.

\q{Afficher le signal carré entre $-10$ et $10$, de période $1$. Utiliser
\texttt{xlim} pour observer un ou deux sauts de plus près.}
\begin{solution}
    \texttt{make\_square(T=1.0, A=1.0)}, puis \texttt{plot\_signal}.
\end{solution}

\begin{lesson}[Fonction \texttt{fourier\_coeffs}]
    \texttt{fourier\_coeffs(signal, N)} calcule $a_0$ et les $N$ premiers
    $a_n,b_n$ par intégration numérique des formules du cours -- pas par
    échantillonnage ni FFT. Retourne \texttt{(a0, an, bn)}.
\end{lesson}

\q{Calculer les 100 premiers coefficients de la série de Fourier du signal
carré. Rappeler la formule des coefficients $a_n$ et $b_n$ vue en cours
(TD1) et vérifier l'accord avec vos résultats.}
\begin{solution}
    $a_n=0$ pour tout $n$ ; $b_n=4/(n\pi)$ si $n$ impair, $b_n=0$ si $n$
    pair. Les valeurs calculées concordent (aux erreurs d'arrondi près sur
    $a_n$, de l'ordre de $10^{-17}$).
\end{solution}

\q{Tracer le spectre d'amplitude correspondant (\texttt{plot\_spectrum}).}
\begin{solution}
    Raies non nulles uniquement aux rangs impairs, décroissant en $1/n$.
\end{solution}

\q{Reconstruire le signal carré à partir de sa série de Fourier pour une
valeur de $N$ de votre choix (\texttt{plot\_reconstruction}), puis
recommencer avec d'autres valeurs (essayer $N=2,5,10,100$). Zoomer avec
\texttt{xlim} près d'un saut. Que remarquez-vous ?}
\begin{solution}
    Plus $N$ augmente, plus la reconstruction se rapproche du signal carré
    loin des sauts. Mais au voisinage de chaque discontinuité, un
    dépassement d'environ 9\% de l'amplitude du saut persiste quel que soit
    $N$ : il se resserre autour du saut sans jamais disparaître. C'est le
    phénomène de Gibbs.
\end{solution}

\q{Reconstruire à nouveau le signal, mais en utilisant seulement les
coefficients de rang \texttt{n\_min} à \texttt{n\_max} (essayer d'abord 50
à 100, puis d'autres bornes). Que remarquez-vous ?}
\begin{solution}
    Le résultat n'a plus l'allure d'un carré : c'est une oscillation rapide
    de faible amplitude (les $b_n$ de rang 50-100 valent
    $4/(n\pi)\approx 0{,}01$--$0{,}025$, bien plus petits que l'amplitude 1
    d'origine). Les harmoniques basses portent la forme globale du signal ;
    les hautes n'ajoutent que du détail fin, avec une énergie bien plus
    faible.
\end{solution}


%% ─── SECTION C ──────────────────────────────────────────────────
\section{Deux signaux à l'allure trompeuse}

Deux signaux fournis vont mettre à l'épreuve votre intuition sur les
conditions de Dirichlet : l'un a l'air pathologique et ne l'est pas,
l'autre l'inverse. On les étudie l'un après l'autre, en entier.

\subsection*{La marche aléatoire}

\q{Afficher le signal \texttt{marche} (\texttt{make\_random\_walk}) sur
quelques périodes. Est-il périodique ? Respecte-t-il les conditions de
Dirichlet ? Justifier.}
\begin{solution}
    Oui, périodique par construction. Il respecte Dirichlet malgré son
    allure erratique : il est construit à partir d'un nombre \emph{fixe} de
    sommets (25) reliés par des segments de droite -- donc un nombre fini
    d'extrema, aucune discontinuité (la dernière valeur est recalée sur la
    première), et borné.
\end{solution}

\q{Calculer ses coefficients de Fourier et reconstruire le signal pour une
valeur de $N$ de votre choix, puis recommencer avec d'autres valeurs. Le
comportement observé est-il cohérent avec la réponse précédente ?}
\begin{solution}
    Oui : la reconstruction converge proprement vers le signal quand $N$
    augmente, cohérent avec Dirichlet respecté.
\end{solution}

\subsection*{Le signal exotique}

\q{Afficher le signal \texttt{exotique} (\texttt{make\_exotique}) sur
quelques périodes, puis zoomer avec \texttt{xlim} près de $t=0$. Est-il
périodique ? Respecte-t-il les conditions de Dirichlet ? Justifier.}
\begin{solution}
    Périodique, mais il oscille infiniment vite au voisinage de $t=0$ dans
    chaque période ($\sin(2\pi/t)$ a une infinité d'extrema qui
    s'accumulent quand $t\to0^+$, bien visible en zoomant). Il ne respecte
    donc \emph{pas} les conditions de Dirichlet (nombre fini d'extrema),
    bien que son tracé à l'échelle globale ait l'air raisonnable.
\end{solution}

\q{Calculer ses coefficients de Fourier et reconstruire le signal pour une
valeur de $N$ de votre choix, puis recommencer avec d'autres valeurs. Le
comportement observé est-il cohérent avec la réponse précédente ?}
\begin{solution}
    Oui : le calcul des coefficients affiche lui-même un avertissement
    (« le calcul numérique peine à converger »), et la reconstruction reste
    mauvaise près de $t=0$ même à $N=60$ -- cohérent avec la
    non-convergence attendue pour un signal qui ne vérifie pas Dirichlet.
\end{solution}


%% ─── SECTION D (optionnelle) ─────────────────────────────────────
\section{Optionnel, si le temps le permet --- le signal dents de scie}

\q{Afficher le signal \texttt{scie} (\texttt{make\_sawtooth}) sur quelques
périodes. Où se trouve sa discontinuité, comparée à celle du signal carré ?}
\begin{solution}
    Une seule discontinuité par période, exactement à la jonction $t=0/T$
    (contrairement au carré, qui en a une au milieu de la période \emph{et}
    une à la jonction).
\end{solution}

\q{Calculer ses coefficients de Fourier et reconstruire le signal pour une
valeur de $N$ de votre choix. Zoomer avec \texttt{xlim} exactement sur la
jonction de période. Que remarquez-vous ?}
\begin{solution}
    Même phénomène que pour le carré : un dépassement d'environ 9\% de
    l'amplitude du saut apparaît juste avant/après la jonction et ne
    disparaît pas quand $N$ augmente -- Gibbs, ici localisé au seul endroit
    où le signal est discontinu.
\end{solution}

\q{Tracer le spectre d'amplitude et le comparer à celui du signal carré
(question B3).}
\begin{solution}
    Ici $b_n=-1/(n\pi)$ pour \emph{tout} $n$ (spectre dense), contre
    $b_n=4/(n\pi)$ sur les rangs \emph{impairs seulement} pour le carré
    (spectre creux, amplitude plus grande par raie). Mais dans les deux cas
    l'enveloppe décroît en $1/n$ : attendu, puisque les deux signaux ont
    une discontinuité de première espèce -- c'est cette décroissance lente
    qui est responsable du phénomène de Gibbs dans les deux cas.
\end{solution}

\begin{lesson}[Fonctions \texttt{parseval\_energy}, \texttt{signal\_energy}]
    \texttt{parseval\_energy(a0, an, bn)} calcule
    $a_0^2+\frac12\sum_n(a_n^2+b_n^2)$ à partir des coefficients seuls.
    \texttt{signal\_energy(signal)} calcule $\frac1T\int_0^T f(t)^2\,\d t$
    directement sur le signal. Deux façons de calculer la même chose ?
\end{lesson}

\q{Calculer \texttt{signal\_energy(scie)}, puis \texttt{parseval\_energy}
pour plusieurs valeurs de $N$ croissantes. Que remarquez-vous ?}
\begin{solution}
    La somme partielle $a_0^2+\frac12\sum(a_n^2+b_n^2)$ se rapproche, quand
    $N$ augmente, de la même valeur que l'intégrale directe
    $\frac1T\int_0^T f(t)^2\,\d t$ (ici $\approx 0{,}333$ dans les deux
    cas). L'« énergie » du signal peut donc se lire soit directement sur
    $f$, soit sur son spectre. On retrouvera ce résultat plus tard dans le
    cours sous le nom de théorème de Parseval.
\end{solution}

\end{document}
